% ═══════════════════════════════════════════════════════════════
% SL-07d-ML: Stundenleistung MUSTERLÖSUNG
% Spanisch Klasse 7 - Sequenz 7: Compras
% ═══════════════════════════════════════════════════════════════

\documentclass[11pt, a4paper]{article}

\newif\ifswmodus
\swmodusfalse

\usepackage[utf8]{inputenc}
\usepackage[T1]{fontenc}
\usepackage[ngerman]{babel}
\usepackage{helvet}
\renewcommand{\familydefault}{\sfdefault}

\usepackage[margin=2cm]{geometry}
\usepackage{enumitem}
\usepackage{tabularx}
\usepackage{booktabs}

\usepackage{xcolor}
\usepackage{tcolorbox}
\tcbuselibrary{skins}

\usepackage{fancyhdr}
\usepackage{lastpage}
\usepackage{needspace}

\definecolor{spanisch-color}{HTML}{c41e3a}
\definecolor{grau-color}{HTML}{888888}
\definecolor{hellgrau-color}{HTML}{f5f5f5}
\definecolor{loesung-color}{HTML}{c62828}

\definecolor{sw-dark}{HTML}{000000}
\definecolor{sw-gray}{HTML}{666666}
\definecolor{sw-white}{HTML}{ffffff}

\ifswmodus
    \colorlet{spanisch}{sw-dark}
    \colorlet{grau}{sw-gray}
    \colorlet{hellgrau}{sw-white}
    \colorlet{loesung}{sw-dark}
\else
    \colorlet{spanisch}{spanisch-color}
    \colorlet{grau}{grau-color}
    \colorlet{hellgrau}{hellgrau-color}
    \colorlet{loesung}{loesung-color}
\fi

\newcommand{\fachfarbe}{spanisch}

\newcommand{\thema}{Stundenleistung: Compras}
\newcommand{\fach}{Spanisch}
\newcommand{\klasse}{7}
\newcommand{\maxpunkte}{20}
\newcommand{\zeit}{25 Minuten}
\newcommand{\datum}{\today}

\pagestyle{fancy}
\fancyhf{}
\renewcommand{\headrulewidth}{0.4pt}
\renewcommand{\footrulewidth}{0.4pt}

\lhead{\fach{} -- Klasse \klasse}
\chead{\textcolor{loesung}{\textbf{SL-07d MUSTERLÖSUNG}}}
\rhead{\datum}

\lfoot{\zeit}
\cfoot{}
\rfoot{Seite \thepage\ von \pageref{LastPage}}

\newcommand{\punktebox}[1]{\hfill\fbox{\small \_/#1}}
\newcommand{\luecke}[1]{\rule{#1}{0.4pt}}
\newcommand{\lsg}[1]{\textcolor{loesung}{\textbf{#1}}}

\newtcolorbox{infobox}{
    colback=hellgrau,
    colframe=\fachfarbe,
    boxrule=1pt,
    arc=0pt,
    left=8pt, right=8pt, top=8pt, bottom=8pt
}

\newenvironment{aufgabe}[1]{%
    \needspace{5\baselineskip}%
    \nopagebreak[4]%
    \vspace{10pt}
    \noindent\textbf{\textcolor{\fachfarbe}{Aufgabe #1}}
    \nopagebreak[4]%
    \vspace{5pt}
    \nopagebreak[4]%
    \begin{list}{}{\leftmargin=0pt}
    \item[]
}{%
    \end{list}
}

\newcommand{\es}[1]{\textcolor{\fachfarbe}{\textbf{#1}}}

\begin{document}

\begin{center}
    {\Large\bfseries\textcolor{\fachfarbe}{\thema}}\\[0.3em]
    {\large\textcolor{loesung}{\textbf{--- MUSTERLÖSUNG ---}}}
\end{center}

\vspace{5pt}

\noindent
\begin{tabularx}{\textwidth}{|X|X|X|}
\hline
\textbf{Name:} \lsg{Musterschüler/in} &
\textbf{Klasse:} \klasse &
\textbf{Datum:} \datum \\
\hline
\end{tabularx}

\vspace{10pt}

\begin{infobox}
\textbf{Zeit:} \zeit{} | \textbf{Punkte:} \maxpunkte{} | \textbf{Hilfsmittel:} keine
\end{infobox}

\vspace{15pt}

\begin{aufgabe}{1 -- Lebensmittel zuordnen}
Verbinde das spanische Wort mit der deutschen Bedeutung. \punktebox{4}
\end{aufgabe}

\begin{center}
\begin{tabular}{rcl}
\es{el queso} & \lsg{---} & das Brot \lsg{(= el queso $\rightarrow$ der Käse)} \\[0.8em]
\es{la manzana} & \lsg{---} & der Käse \lsg{(= la manzana $\rightarrow$ der Apfel)} \\[0.8em]
\es{el pan} & \lsg{---} & das Hähnchen \lsg{(= el pan $\rightarrow$ das Brot)} \\[0.8em]
\es{el pollo} & \lsg{---} & der Apfel \lsg{(= el pollo $\rightarrow$ das Hähnchen)} \\
\end{tabular}
\end{center}

\textbf{Bewertung:} Je 1 Punkt pro richtige Zuordnung (4P)

\vspace{15pt}

\begin{aufgabe}{2 -- Zahlen als Wort schreiben}
Schreibe die Zahl als spanisches Wort. \punktebox{4}
\end{aufgabe}

\noindent
a) 15 = \lsg{quince}

\vspace{0.8em}

\noindent
b) 32 = \lsg{treinta y dos}

\vspace{0.8em}

\noindent
c) 67 = \lsg{sesenta y siete}

\vspace{0.8em}

\noindent
d) 100 = \lsg{cien}

\textbf{Bewertung:} Je 1 Punkt pro richtige Schreibweise (4P)

\vspace{15pt}

\begin{aufgabe}{3 -- Konjugiere \es{querer} und \es{poder}}
Ergänze die richtige Form. \punktebox{6}
\end{aufgabe}

\noindent
a) Yo \lsg{quiero} comprar naranjas. \hfill (querer -- wollen)

\vspace{0.8em}

\noindent
b) ¿Tú \lsg{puedes} ayudarme? \hfill (poder -- können)

\vspace{0.8em}

\noindent
c) María \lsg{quiere} un kilo de tomates. \hfill (querer -- wollen)

\vspace{0.8em}

\noindent
d) Nosotros no \lsg{podemos} pagar con tarjeta. \hfill (poder -- können)

\vspace{0.8em}

\noindent
e) Ellos \lsg{quieren} ir al mercado. \hfill (querer -- wollen)

\vspace{0.8em}

\noindent
f) ¿Vosotros \lsg{podéis} hablar español? \hfill (poder -- können)

\textbf{Bewertung:} Je 1 Punkt pro richtige Form (6P)

\vspace{15pt}

\begin{aufgabe}{4 -- Einkaufsdialog vervollständigen}
Ergänze den Dialog sinnvoll auf Spanisch. \punktebox{6}
\end{aufgabe}

\noindent
\textbf{Verkäufer:} ¡Buenos días! ¿Qué desea?

\vspace{0.8em}

\noindent
\textbf{Kunde:} \lsg{Quiero un kilo de naranjas.} \hfill (2P)

\vspace{0.8em}

\noindent
\textbf{Verkäufer:} Aquí tiene. ¿Algo más?

\vspace{0.8em}

\noindent
\textbf{Kunde:} Sí, \lsg{¿puedo tener queso?} \hfill (1P)

\vspace{0.8em}

\noindent
\textbf{Verkäufer:} Claro. ¿Cuánto quiere?

\vspace{0.8em}

\noindent
\textbf{Kunde:} 200 gramos. \lsg{¿Cuánto es?} \hfill (1P)

\vspace{0.8em}

\noindent
\textbf{Verkäufer:} \lsg{Son seis euros cincuenta.} \hfill (2P)

\textbf{Bewertung:}
\begin{itemize}[noitemsep]
    \item Kunde 1: 2P (Quiero + Menge + Lebensmittel)
    \item Kunde 2: 1P (puedo tener / quisiera + Lebensmittel)
    \item Kunde 3: 1P (¿Cuánto es? / ¿Cuánto cuesta?)
    \item Verkäufer: 2P (Son + korrekte Preisangabe)
\end{itemize}

\vspace{25pt}
\hrule
\vspace{10pt}

\noindent
\begin{tabularx}{\textwidth}{|X|c|c|c|c||c|}
\hline
\textbf{Aufgabe} & 1 & 2 & 3 & 4 & \textbf{Gesamt} \\
\hline
\textbf{max. Punkte} & 4 & 4 & 6 & 6 & \textbf{20} \\
\hline
\textbf{erreicht} & \lsg{4} & \lsg{4} & \lsg{6} & \lsg{6} & \lsg{20} \\
\hline
\end{tabularx}

\vspace{10pt}

\begin{flushright}
\textbf{Note:} \lsg{1}
\end{flushright}

\vspace{1em}

\textbf{Notenschlüssel (Klasse 7):}

\begin{tabular}{|c|c|c|c|c|c|c|}
\hline
\textbf{Note} & 1 & 2 & 3 & 4 & 5 & 6 \\
\hline
\textbf{ab \%} & 96 & 80 & 60 & 40 & 20 & <20 \\
\hline
\textbf{ab Punkte} & 19 & 16 & 12 & 8 & 4 & <4 \\
\hline
\end{tabular}

\vspace{1em}

\textbf{Hinweise zur Korrektur:}
\begin{itemize}[noitemsep]
    \item Aufgabe 1: Linien müssen eindeutig die richtigen Paare verbinden
    \item Aufgabe 2: Rechtschreibfehler bei y/i tolerieren (z.B. ``treinta i dos'')
    \item Aufgabe 3: Stammvokalwechsel muss korrekt sein (e$\rightarrow$ie, o$\rightarrow$ue)
    \item Aufgabe 4: Sinngemäße Antworten akzeptieren (z.B. ``Quisiera'' statt ``Quiero'')
\end{itemize}

\end{document}
